\section{A thermal block entropy inequality}\label{sec:thermal}

The zero-block theorem does not apply directly when thermal excitation
and decay both occur. We will express the resulting state as a convex
mixture of two states to which it does apply, preserving the same scalar
entropy term in each. We first recall two elementary entropy facts,
including their proofs to fix the directions of the inequalities used
below. For a state $P$,
\begin{equation}\label{eq:diagonal}
 \ent(P)\le H(\diag P).
\end{equation}
Indeed, its diagonal is obtained from its eigenvalue vector by the
doubly stochastic matrix of squared absolute unitary entries.
Concavity of $-s\ln s$, applied row by row and then summed, gives
\eqref{eq:diagonal}. Entropy is also concave on states:
\begin{equation}\label{eq:concavity}
 \ent\!\left(\sum_iw_iP_i\right)\ge\sum_iw_i\ent(P_i).
\end{equation}
To see this, choose a basis diagonalizing the mixture, apply classical
entropy concavity in that basis, and then use \eqref{eq:diagonal} on
each $P_i$.

The next result is the extension needed for thermal channels.

\begin{theorem}\label{thm:thermal}
Let $X,Y$ be complex matrices of the same rectangular shape with
$\Tr(XX^*)=\Tr(YY^*)=1$. Suppose
$\alpha,\beta,\eta,\delta\ge0$ sum to one and
\begin{equation}\label{eq:thermal-condition}
 \alpha\delta-\beta\eta=K^2,\qquad K\ge0.
\end{equation}
Set
\begin{equation}\label{eq:thermal-state}
 T=\begin{pmatrix}
 \alpha XX^*+\beta YY^*&KXY^*\\
 KYX^*&\eta XX^*+\delta YY^*
 \end{pmatrix},\qquad
 U_0=\alpha+\beta,\quad U_1=\eta+\delta,
\end{equation}
and $u=U_0U_1-K^2$. Then $T$ is a state, $0\le u\le1/4$, and
\begin{equation}\label{eq:thermal-bound}
 \ent(T)\ge(\alpha+\eta)\ent(XX^*)
             +(\beta+\delta)\ent(YY^*)+g(u).
\end{equation}
\end{theorem}

\begin{proof}
The coefficient identity gives
\begin{equation}\label{eq:thermal-u}
 u=\beta U_1+\eta U_0.
\end{equation}
It follows that $0\le u\le U_0U_1\le1/4$.

Suppose first that $u>0$, so $U_0,U_1>0$. To construct the mixture,
keep $X,Y,K$ and the row sums $U_0,U_1$ fixed. Writing
$\alpha=U_0-\beta$ and $\delta=U_1-\eta$, the coefficient condition
\eqref{eq:thermal-condition} becomes the line equation
$U_1\beta+U_0\eta=u$. Its nonnegative solutions form the segment
joining $(\beta,\eta)=(u/U_1,0)$ and $(0,u/U_0)$. At either endpoint
one coefficient vanishes, permitting a zero-block Gram factorization.
The fixed row sums and cross coefficient retain the same value of $u$.
These endpoints and the weights expressing the original coefficients
as their convex mixture are
\begin{equation}\label{eq:thermal-table}
 \begin{array}{c|cccc|c}
 i&\alpha_i&\beta_i&\eta_i&\delta_i&w_i\\ \hline
 1&K^2/U_1&u/U_1&0&U_1&\beta U_1/u\\[2pt]
 2&U_0&0&u/U_0&K^2/U_0&\eta U_0/u
 \end{array}
\end{equation}
Both coefficient sets are nonnegative and have the original row sums
$U_0,U_1$, since $K^2+u=U_0U_1$. Each therefore sums to one and
satisfies $\alpha_i\delta_i=K^2$. Moreover, $w_1+w_2=1$ by
\eqref{eq:thermal-u}, and
\[
 \sum_iw_i\beta_i=\beta,\qquad
 \sum_iw_i\eta_i=\eta.
\]
The fixed row sums then recover $\alpha$ and $\delta$ as well.

Let $T_i$ denote \eqref{eq:thermal-state} with the $i$th coefficient
set. Both $T_i$ have cross block $KXY^*$, so
$T=w_1T_1+w_2T_2$. They have Gram factorizations
\begin{equation}\label{eq:thermal-factors}
 \begin{split}
 M_1&=\begin{pmatrix}
 \sqrt{\alpha_1}X&\sqrt{\beta_1}Y\\
 \sqrt{\delta_1}Y&0
 \end{pmatrix},\\
 M_2&=\begin{pmatrix}
 \sqrt{\alpha_2}X&0\\
 \sqrt{\delta_2}Y&\sqrt{\eta_2}X
 \end{pmatrix},\qquad T_i=M_iM_i^*.
 \end{split}
\end{equation}
In particular, $T_i$ and $T$ are positive semidefinite of trace one.
Apply \Cref{thm:block} to $M_1$ and to $M_2$ after swapping its two
block rows. The swap is unitary on the output space and preserves
entropy. The respective products of the off-corner block masses are
\[
 \beta_1\delta_1=u,\qquad \eta_2\alpha_2=u.
\]
Thus \emph{the same} scalar term occurs in both inequalities:
\[
 \ent(T_i)\ge(\alpha_i+\eta_i)\ent(XX^*)
                 +(\beta_i+\delta_i)\ent(YY^*)+g(u).
\]
Concavity \eqref{eq:concavity} and recovery of the original
coefficients now give \eqref{eq:thermal-bound}. Zero weights in
\eqref{eq:thermal-table} do not change this argument.

It remains to consider $u=0$. If a row sum vanishes, all its
nonnegative coefficients vanish and $K=0$. The other diagonal block
is a convex mixture of $XX^*$ and $YY^*$, so positivity, normalization,
and \eqref{eq:thermal-bound} follow from \eqref{eq:concavity} and
$g(0)=0$. If both row sums are positive, \eqref{eq:thermal-u} forces
$\beta=\eta=0$, and $K=\sqrt{\alpha\delta}$. Then the first
factorization in \eqref{eq:thermal-factors}, with coefficients
$(\alpha,0,0,\delta)$, applies directly. It proves positivity and
reduces \eqref{eq:thermal-bound} to \Cref{thm:block} with a zero block
mass. These cases exhaust the possibilities.
\end{proof}

\begin{remark}\label{rem:thermal-sharp}
No orthogonality condition on $X,Y$ and no commutativity of their Gram
products is used. The bound is sharp when $X=Y$: in that case
$T=\bigl(\begin{smallmatrix}U_0&K\\K&U_1\end{smallmatrix}\bigr)
\otimes XX^*$, whose entropy is $g(u)+\ent(XX^*)$.
\end{remark}
